\chapter*{序言}
\addcontentsline{toc}{chapter}{序言}

这本书是整个系列的第一册。它的目标不是停留在“看懂一些汇编”，而是建立一套能继续进入系统、编译器、体系结构和性能工程的基础模型：

\begin{center}
\begin{minipage}{0.82\linewidth}
\begin{verbatim}
C/C++ 源码语义
  -> 编译、汇编、链接
  -> x86-64 汇编
  -> 指令、寄存器和内存
  -> ABI、调用约定和二进制接口
  -> 工具观察和可信性能测量
\end{verbatim}
\end{minipage}
\end{center}

本书面向已经写过基本 C/C++ 程序、希望进入系统、编译器和性能工程的人。正文直接从程序到进程开始，再进入优化器、CPU 执行、对象布局、工具、x86-64、ABI、互操作和可信测量。

本书不是指令速查手册。每个结论都要能落到源码、二进制、机器状态或实验输出上：只看源码不够，只看汇编不够，只看运行时间也不够。读者需要运行命令、观察输出、写代码、看反汇编、做 benchmark，并保存能复查的报告。
